% !PRINTED 17
\weberpage{17}

%<<CONT
wo $H$ und $T$ Formen vierten und sechsten Grades sind, deren
Koeffizienten sich rational und mit ganzzahligen Zahlenkoeffizienten
aus den Koeffizienten von $f$ zusammensetzen. Zwischen diesen
Formen besteht dann noch die Relation
\weq{15}{H^3 - 48\,AHf^2 - 64\,Bf^3 = -27\,T^2}
(Bd. I, \S\ 70, 72).

Es ist aber nach dem \gesperrt{Eulerschen} Satz über homogene
Funktionen
\[
\begin{aligned}
4f&=\frac{\partial f}{\partial x}x+\frac{\partial f}{\partial y}y,
&df&=\frac{\partial f}{\partial x}dx+\frac{\partial f}{\partial y}dy,\\
4H&=\frac{\partial H}{\partial x}x+\frac{\partial H}{\partial y}y,
&dH&=\frac{\partial H}{\partial x}dx+\frac{\partial H}{\partial y}dy,
\end{aligned}
\]
woraus
\[
f\,dH-H\,df=-3\,T(y\,dx-x\,dy),
\]
und wenn also nach (15)
\weq{16}{3\sqrt{3}\,T = \sqrt{-H^3+48\,Af^2H+64\,Bf^3}}
gesetzt wird:
\[
\frac{y\,dx-x\,dy}{\sqrt{f}} = \frac{-\sqrt{3}\,(f\,dH-H\,df)}{\sqrt{-(H^3-48\,Af^2H-64\,Bf^3)f}}.
\]

Macht man nun die Substitution
\weq{17}{z = -\frac{\mu H}{4f}}
mit einem unbestimmten Faktor $\mu$, so ergibt sich
\weq{18}{\frac{y\,dx-x\,dy}{\sqrt{3\mu f}} = \frac{dz}{\sqrt{4z^3-g_2z-g_3}},}
wenn wie früher
\weq{19}{g_2 = 12\,\mu^2A,\qquad g_3=-4\,\mu^3B.}

Man kann diese Transformation auf ein Differential anwenden,
das schon die Normalform hat. Setzt man
\[
f(x,y) = 4x^3y - g_2xy^3 - g_3y^4,
\]
so hat man
\[
a_0\quad a_1\quad a_2\quad a_3\quad a_4
\]
durch
\[
0\quad 4\quad 0\quad -g_2\quad -g_3
\]
zu ersetzen, und es ergibt sich aus (7)
\[
A = 12\,g_2, \qquad B = -27\cdot 16\,g_3.
\]
